\chapter*{Resumen}

\section*{\tituloPortadaVal}

El progreso en un deporte individual como la escalada suele ser lento y estar marcado por largas fases de estancamiento, que suponen un reto importante para mantener la motivación y la constancia del escalador. A este problema se suma que el seguimiento del progreso y los logros se realizan de forma manual en la mayoría de los casos, recurriendo a la memoria o libretas. 

El principal objetivo de este trabajo consiste en aliviar estos dos problemas mediante una aplicación para dispositivos móviles que digitaliza la gestión de las rutas en los rocódromos y facilita a los escaladores un seguimiento exhaustivo de su actividad. Para atajar el problema de la desmotivación, la herramienta incorpora mecánicas de gamificación basadas en el sistema PBL (puntos, insignias y clasificaciones) cuyos elementos transforman la rutina de entrenamiento en una experiencia más dinámica, permitiendo visualizar el progreso a los usuarios e incentivando una competitividad sana que refuerza el sentido de pertenencia a la comunidad de escalada. 

El proyecto se ha desarrollado bajo un enfoque DCU (Diseño Centrado en Usuario), logrando de este modo que la definición de requisitos, las propuestas de mejora y las pruebas realizadas surgen directamente de las necesidades de los escaladores. De este modo la solución construida aporta un valor real y se adapta a las necesidades de los usuarios finales.

\section*{Palabras clave}
   
\noindent Escalada indoor, gamificación, rocódromos, motivación, sistema PBL, diseño centrado en el usuario, aplicación móvil, seguimiento deportivo.